% Part: lambda-calculus
% Chapter: lambda-definability
% Section: primitive-recursive-functions

\documentclass[../../../include/open-logic-section]{subfiles}

\begin{document}

\olfileid{lam}{ldf}{prf}
\olsection{આદિમ પુનરાવર્તી વિધેયો \usetoken{S}{lambda definable} છે}

યાદ કરો કે આદિમ પુનરાવર્તી વિધેયો એ છે, જે મૂળભૂત
વિધેયો $\Zero$, $\Succ$ અને $\Proj{n}{i}$માંથી
સંયોજન તથા આદિમ પુનરાવર્તન દ્વારા વ્યાખ્યાયિત કરી શકાય.

\begin{lem}
  \ollabel{lem:basic}
  મૂળભૂત આદિમ પુનરાવર્તી વિધેયો $\Zero$, $\Succ$ અને
  પ્રક્ષેપો~$\Proj{n}{i}$ !!{lambda definable} છે.
\end{lem}

\begin{proof}
નીચેનાં પદો તેમને !!{lambda define} કરે છે:
\begin{align*}
  \fn{Zero} & \ident \lambd[a][\lambd[fx][x]]\\
  \fn{Succ} & \ident \lambd[a][\lambd[fx][f (a f x)]] \\
  \fn{Proj}^n_i & \ident \lambd[x_0\dots x_{n-1}][x_i]
\end{align*}
\end{proof}

\begin{lem}
  \ollabel{lem:comp} ધારો કે $k$-સ્થાનીય વિધેય $f$ અને
  $n$-સ્થાનીય વિધેયો $g_0, \dots, g_{k-1}$ અનુક્રમે
  $F$, $G_0$, \dots, $G_{k-1}$ પદો દ્વારા
  !!{lambda definable} છે, અને $h$ તેમની
  સંયોજનથી વ્યાખ્યાયિત થાય છે. તો $h$ પણ
  !!{lambda definable} છે.
  \sourcecorrection{OLLAM-054}{સ્થિર અંગ્રેજી મૂળમાં
  $k$ વિધેયો $g_0,\dots,g_{k-1}$ હોવા છતાં
  તેમને દર્શાવતાં પદોનો ક્રમ $G_k$ સુધી જાય છે;
  અહીં છેલ્લું પદ $G_{k-1}$ કર્યું છે. મૂળના
  નિષ્કર્ષમાં ``$H$ is lambda definable'' છે,
  પરંતુ $H$ નિરૂપણ પદ છે; વ્યાખ્યાયિત થતું
  વિધેય $h$ છે.}
\end{lem}

\begin{proof}
  વિધેય~$h$ને આ પદ !!{lambda define} કરે છે:
  \[
  H \ident \lambd[x_0 \dots x_{n-1}][F\, (G_0 x_0 \dots
    x_{n-1}) \dots (G_{k-1} x_0 \dots x_{n-1})]
  \]
  આ હકીકતની ચકાસણી કસરત તરીકે છોડી છે.
\end{proof}

\begin{prob}
  $H\num{n_0}\dots\num{n_{n-1}} \red \num{h(n_0, \dots,
    n_{n-1})}$ બતાવીને
  \olref[lam][ldf][prf]{lem:comp}નો પુરાવો પૂર્ણ કરો.
\end{prob}

નોંધો કે \olref{lem:comp}માં $f$ અને $g_0$,
\dots,~$g_{k-1}$ આદિમ પુનરાવર્તી હોવાં જરૂરી
નથી; તેઓ સર્વત્ર વ્યાખ્યાયિત અને !!{lambda definable}
હોય એટલું જ જરૂરી છે.

\begin{lem}\ollabel{lem:prim}
  ધારો કે $f$ એ $n$-સ્થાનીય વિધેય અને~$g$ એ
  $n+2$-સ્થાનીય વિધેય છે; તેઓ અનુક્રમે
  $F$ અને~$G$ પદો દ્વારા !!{lambda definable} છે,
  અને $h$ એ $f$ તથા $g$માંથી આદિમ
  પુનરાવર્તન દ્વારા વ્યાખ્યાયિત છે. તો $h$ પણ
  !!{lambda definable} છે.
\end{lem}

\begin{proof}
  યાદ કરો કે $h$ની વ્યાખ્યા આ છે:
  \begin{align*}
    h(x_1, \dots, x_n, 0) &= f(x_1, \dots, x_n)\\
    h(x_1, \dots, x_n, y+1) & = g(x_1, \dots, x_n, y, h(x_1, \dots, x_n, y)).
  \end{align*}
  \sourcecorrection{OLLAM-055}{સ્થિર અંગ્રેજી મૂળના
  બીજા પુનરાવર્તન સમીકરણમાં જમણી બાજુનું બહારનું
  વિધેય ભૂલથી $h$ છે. પગલું-વિધેય $g$ હોવું
  જોઈએ; પછીના પુરાવામાં પણ $g$ જ વપરાય છે.}
  અનૌપચારિક રીતે, આદિમ પુનરાવર્તનની વ્યાખ્યા
  $f(x_1,\dots,x_n)$થી શરૂ કરીને પગલું-વિધેય~$g$
  $y$ વાર લાગુ કરે છે. આ ચર્ચ અંકોની પુનરાવર્તક
  તરીકેની વ્યાખ્યાની યાદ અપાવે છે.

  સરળતા માટે, એક જ વધારાની દલીલ~$x$ હોય તે
  કિસ્સાની વ્યાખ્યા અને સાબિતી આપીએ છીએ.
  વિધેય~$h$ને આ પદ !!{lambda define} કરે છે:
  \begin{align*}
    H \ident & \lambd[x][\lambd[y][\fn{Snd} (y
        D \tuple{\num{0}, F x})]]
    \intertext{જ્યાં }
    D \ident & \lambd[p][\tuple{\fn{Succ} (\fn{Fst}\, p), (G x
          (\fn{Fst}\, p) (\fn{Snd}\, p))}]
  \end{align*}
  \sourcecorrection{OLLAM-056}{અહીં $D$માં $x$
  મુક્ત રહે છે; $H$ના બહારના $\lambd[x]$ની અંદર
  તેનો અર્થ દરેક નક્કી કરેલા $x$ માટેનું પગલું-પદ
  છે. પછીના પુરાવાનું $D_n$ એ $x$ના સ્થાને
  $\num n$ મૂકેલું આ જ પદ છે. આ અવલંબન
  મૂળના સંકેતનમાં સ્પષ્ટ નહોતું.}
  આપણે રાખેલી પુનરાવર્તનની સ્થિતિ એક જોડી છે:
  તેનો પહેલો ઘટક હાલનું~$y$ છે અને બીજો~$h$નું
  અનુરૂપ મૂલ્ય છે. દરેક પગલે $y$ અને~$h$નાં
  નવા મૂલ્યોની જોડી બનાવીએ છીએ; પુનરાવર્તન પૂરું
  થયા પછી તેનો બીજો ઘટક પાછો આપીએ છીએ;
  તે જ $h$નું અંતિમ મૂલ્ય છે.
  હવે આ ખરેખર આદિમ પુનરાવર્તનનું નિરૂપણ છે
  તે સાબિત કરીએ.

  આપણે દરેક $n$ અને $m$ માટે
  $H\,\num{n}\,\num{m} \red \num{h(n,m)}$
  સાબિત કરવું છે. પહેલાં, $D_n \ident
  \Subst{D}{\num{n}}{x}$ હોય તો
  $D_n^m \tuple{\num{0}, F\, \num{n}} \red
  \tuple{\num{m}, \num{h(n, m)}}$ છે એમ
  $m$ પર ગાણિતિક અનુમાનથી બતાવીએ.

  $m=0$ હોય તો
  $D_n^0 \tuple{\num{0}, F\, \num{n}} \red
  \tuple{\num{0}, \num{h(n, 0)}}$ બતાવવું છે.
  પરંતુ $D_n^0 \tuple{\num{0}, F\, \num{n}}$
  એટલે $\tuple{\num{0}, F\, \num{n}}$ જ.
  $F$ એ~$f$ને !!{lambda define} કરતું હોવાથી
  આ $\tuple{\num{0}, \num{f(n)}}$માં
  લઘુકૃત થાય છે; અને $f(n)=h(n,0)$ હોવાથી
  તે $\tuple{\num{0}, \num{h(n,0)}}$ છે.

  હવે ધારો કે
  $D_n^m \tuple{\num{0}, F\, \num{n}} \red
  \tuple{\num{m}, \num{h(n, m)}}$.
  આપણે $D_n^{m+1} \tuple{\num{0}, F\, \num{n}} \red
  \tuple{\num{m+1}, \num{h(n, m+1)}}$
  બતાવવું છે.
  \begin{align*}
    D_n^{m+1} \tuple{\num{0}, F\, \num{n}}
    & \ident D_n(D_n^{m} \tuple{\num{0}, F\, \num{n}})\\
    & \red D_n\,\tuple{\num{m}, \num{h(n, m)}} \text{\qquad (અનુમાનથી)}\\
    & \ident (\lambd[p][
      \tuple{
        \fn{Succ} (\fn{Fst}\, p), (G \, \num{n} (\fn{Fst}\, p) (\fn{Snd}\, p))
    }]) \tuple{\num{m}, \num{h(n,m)}}\\
    & \redone
    \tuple{\fn{Succ} (\fn{Fst}\, \tuple{\num{m}, \num{h(n,m)}}),\\
      & \qquad (G \, \num{n} (\fn{Fst}\, \tuple{\num{m}, \num{h(n,m)}}) (\fn{Snd}\, \tuple{\num{m}, \num{h(n,m)}}))}\\
    & \red
    \tuple{\fn{Succ}\, \num{m}, (G \, \num{n} \,\num{m} \, \num{h(n,m)})}\\
    & \red \tuple{\num{m+1}, \num{g(n, m, h(n, m))}}
  \end{align*}
  $g(n,m,h(n,m))=h(n,m+1)$ હોવાથી ઇચ્છિત
  પરિણામ મળે છે.

  અંતે, આ ગણતરી જુઓ:
  \begin{align*}
    H\,\num n\,\num m &\ident \lambd[x][\lambd[y][\fn{Snd} (y
        (\lambd[p].\tuple{\fn{Succ} (\fn{Fst}\, p), (G\, x\,
          (\fn{Fst}\, p)\, (\fn{Snd}\, p))}) \tuple{\num{0}, F x})]]\\
    & \qquad\,\num n\, \num m\\
    & \red \fn{Snd} (\num m \,
    \underbrace{(\lambd[p].\tuple{\fn{Succ} (\fn{Fst}\, p), (G \,\num n\,
      (\fn{Fst}\, p) (\fn{Snd}\, p))})}_{D_n} \tuple{\num{0}, F \num n})\\
    & \ident \fn{Snd} (\num m \, D_n\, \tuple{\num{0}, F \num n})\\
    & \red \fn{Snd} \, (D_n^m \tuple{\num{0}, F \num n})\\
    & \red \fn{Snd} \,\tuple{\num{m}, \num{h(n, m)}} \\
    & \red \num{h(n,m)}.
  \end{align*}
\end{proof}

\begin{prop}
  દરેક આદિમ પુનરાવર્તી વિધેય !!{lambda definable} છે.
\end{prop}

\begin{proof}
  \olref{lem:basic} મુજબ બધાં મૂળભૂત વિધેયો
  !!{lambda definable} છે. \olref{lem:comp}
  અને \olref{lem:prim} મુજબ
  !!{lambda definable} વિધેયો સંયોજન તથા
  આદિમ પુનરાવર્તન હેઠળ બંધ છે.
\end{proof}

\end{document}
